\exnumber{04}
\extitle{UX III - Ex 5 - Personas}
\exfiles{ex04/fichier Figma et PDF, nomme Name\_FirstName\_DeliverableName\_Date\_VersionNumber}
\exauthorize{Those named in the exercise}
\makeheaderfiles

\textbf{Instructions :}

Let's create a persona based on your quantitative research and your interviews.

Your persona card should include common trends amongst the people you got data from (habits, verbatims, fears, drivers, frustration, and pain points).

Reminder: Your persona is fictional and should not have the exact name or picture as someone you interviewed, but remember a persona is here so you can make decisions as a designer.

\textbf{Livrables :}

\begin{itemize}
\item A PDF persona card extracted from a Figma file
\item Déposez le PDF sur Github, nommé ainsi : « Name\_FirstName\_DeliverableName\_Date\_VersionNumber ». Exemple : « Doe\_John\_MidFiPrototype\_05242024\_V1 ».
\item Ajoutez un titre dans le document.
\item Soignez l'aspect global du document. Il doit être clair, simple et lisible. Vous pouvez regarder canva.com, sans surcharger le dessin de détails.
\end{itemize}

\textbf{Ressources :}

\begin{itemize}
\item Make your persona great again \url{https://uxdesign.cc/personas-e60c1c06ead1}
\item Persona, journey maps, user flow, site map, OH MY! \url{https://uxdesign.cc/personas-journey-maps-site-maps-and-user-flows-oh-my-e71d044b4bcb}
\item Persona framework \url{https://miro.medium.com/max/1400/0*RbHfR6BHhhFIENNf}
\end{itemize}

\textbf{Citations :}

\begin{itemize}
\item If you design for everyone, you delight no one.
\item A good user persona is the one based on user research, without regard to how many attributes we can describe.
\end{itemize}

